
puissance. Par exemple sur le topos @@M0@@ annelé par le corps @@M1@@, la représentation @@M2@@ définie par
@@M3@@
n'est pas une puissance. En particulier 6.2.3 ne reste pas valable pour n'importe quel topos localement annelé @@M4@@ et pour n'importe quelle représentation @@M5@@.

\textbf{6.5.3.} En revanche 6.2.3 reste valable pour les représentations associées aux formules (6.2.8), (6.2.9) et (6.2.10) sur un topos annelé quelconque @@M6@@, parce que ces représentations sont «~algébriques~» dans le sens suivant :

Soient @@M7@@ l'objet final de @@M8@@, @@M9@@, et considérons le foncteur
@@M10@@
qui associe à chaque @@M11@@ le schéma affine @@M12@@ sur @@M13@@. Alors pour tout entier positif @@M14@@ le préfaisceau @@M15@@ sur @@M16@@ est composé avec @@M17@@ du préfaisceau correspondant @@M18@@ sur @@M19@@. On dit qu'une représentation sur @@M20@@
@@M21@@
est \emph{algébrique}, si elle provient par @@M22@@ d'une représentation
@@M23@@
de préfaisceaux en groupes sur @@M24@@, i.e., si l'on a
@@M25@@
pour tout @@M26@@, @@M27@@.

La représentation @@M28@@ dans (6.5.4) est bien entendu complètement déterminée par @@M29@@. Il s'ensuit immédiatement que (6.2.3) reste valable dans un topos annelé quelconque, lorsque @@M30@@ est algébrique. Donc en particulier les formules (6.2.8), (6.2.9) et (6.2.10) restent valables sur un topos annelé quelconque.

